\subsection{Synchronous and asynchronous execution}

In synchronous execution, all rollout workers finish a batch before the learner updates.
The behavior and current policies match during generation. They diverge only if the
learner performs multiple optimizer steps on that buffer. After optimization, workers
receive the new policy and generate the next batch.

Here, on-policy data means that the sampled actions come from the policy version serving
as the behavior policy for the current update. Data become off-policy relative to the
learner when a different policy version generated them. The probability ratio measures
that version mismatch; clipping limits how much the learner trusts it.

Asynchronous execution overlaps rollout generation with learning. Suppose a worker loads
policy version 10 and begins a long tool-using trajectory. While it generates, the learner
advances to version 13 using other workers' data. When the trajectory arrives, its
behavior policy remains $\pi_{10}$ while the learner evaluates it under $\pi_{13}$:
\begin{equation}
  \rho_{i,t}
  =\frac{\pi_{13}(o_{i,t}\mid q,o_{i,<t})}
         {\pi_{10}(o_{i,t}\mid q,o_{i,<t})}.
  \label{eq:grpo-asynchronous-ratio}
\end{equation}
This is the same importance-weighting derivation as
Equation~\ref{eq:grpo-trajectory-importance-expectation}. Synchronous reuse creates policy
lag after learner updates; asynchronous execution can create lag before data reaches the
learner.

Importance weighting does not make arbitrarily stale trajectories equivalent to fresh
on-policy data. Exact correction of the prefix distribution would require the high-variance
trajectory product in Equation~\ref{eq:grpo-complete-trajectory-ratio}. The local token
surrogate still trains on prefixes produced by the older worker. Large lag creates extreme
ratios, extensive clipping, and little reliable improvement signal. Asynchronous systems
therefore combine ratios with policy-version metadata, regular worker refreshes, bounded
queues, lag thresholds, and filtering or rejection of stale trajectories. The exact
sampled tokens and their behavior log-probabilities must remain attached to the data;
recomputing the denominator with the learner invalidates the ratio.

This short-lived rollout buffer differs from the Deep Q-Network (DQN) algorithm's
long-lived experience replay \citep{mnih2015dqn}. DQN retains transitions from many policy versions for broad
off-policy reuse and randomly samples them to reduce temporal correlation. GRPO normally
keeps only current or recent behavior-policy data and discards it after limited
optimization. The shared idea is reuse of collected experience; the acceptable policy lag
and the learning objective are different.
